\section{问题四模型的建立与求解}

\subsection{问题四的描述与设计变量}

问题四承接问题三的室外工况，以问题三所得最优坚持时间$t_3^*$作为性能基准，但不再追加材料资金，而是通过改进中间功能层材料，使其在各温度下的放热能力按同一比例提高。由于本问不增加材料用量，三层厚度恢复并固定为原始设计
\begin{equation}
d_1=0.7\ \mathrm{mm},\qquad
d_2=0.4\ \mathrm{mm},\qquad
d_3=0.3\ \mathrm{mm},
\label{eq:q4-base-thickness}
\end{equation}
服装面积、质量和材料费用均不随放热能力提高比例改变。

令$\alpha\ge1$表示新功能材料相对于原材料的放热能力倍率，则所求提高比例为
\begin{equation}
r=\alpha-1,
\end{equation}
最终以$100r\%$报告。问题四的核心任务即为寻找满足坚持时间要求的最小$\alpha$。

\subsection{DSC放热曲线的同比例增强}

问题一已经将附件DSC曲线由持续功率源修正为有限相变过程。对基线校正后的原材料放热峰记为$p(T)\ge0$，在固化温区$[T_f,T_s]$上定义归一化相变权重
\begin{equation}
w(T)=\frac{p(T)}{\displaystyle\int_{T_f}^{T_s}p(\theta)\,\mathrm d\theta},
\qquad
\int_{T_f}^{T_s}w(T)\,\mathrm dT=1.
\label{eq:q4-phase-weight}
\end{equation}
对应相对固化进度为
\begin{equation}
\xi_{\rm DSC}(T)=
\begin{cases}
0,&T\ge T_s,\\[1mm]
\displaystyle\int_T^{T_s}w(\theta)\,\mathrm d\theta,&T_f<T<T_s,\\[2mm]
1,&T\le T_f.
\end{cases}
\label{eq:q4-xi}
\end{equation}

题目规定"各个温度下同比例增加放热能力"，因此新材料的DSC放热曲线写为
\begin{equation}
p_{\alpha}(T)=\alpha p(T).
\label{eq:q4-dsc-scale}
\end{equation}
由于分子、分母同时乘以$\alpha$，式\eqref{eq:q4-phase-weight}中的归一化权重$w(T)$及式\eqref{eq:q4-xi}中的相对固化进度形状均保持不变；改变的是相同固化温区内可释放的总能量尺度。若原功能层单位质量总固化潜热为$L_{\rm pcm}$，则新材料满足
\begin{equation}
L_{\rm pcm}^{(\alpha)}=\alpha L_{\rm pcm}.
\label{eq:q4-latent-scale}
\end{equation}
因此问题一的有效比热修正为
\begin{equation}
\boxed{
c_{2,\rm eff}^{(\alpha)}(T)
=c_2+\alpha L_{\rm pcm}w(T).
}
\label{eq:q4-ceff}
\end{equation}
该表达同时满足两项要求：一是各温度处的相变放热贡献按同一倍率增强；二是相变总能量仍然有限，固化进度达到$\xi=1$后不再继续放热。数值实现继续采用问题一的历史约束，保证固化进度只增不减。

\subsection{给定放热倍率下的传热正演}

问题四保持问题一的环境条件：环境温度$T_\infty=-40\,^{\circ}\mathrm C$、无风、实验者静止站立。人体按附件测试条件作为$37\,^{\circ}\mathrm C$恒温热源，人体侧对流系数仍取$h_b=3.0\ \mathrm{W/(m^2\cdot K)}$。因此给定$\alpha$后，三层控制方程为
\begin{equation}
\left\{
\begin{aligned}
\rho_1c_1T_{1,t}&=k_1T_{1,xx},\\
\rho_2c_{2,\rm eff}^{(\alpha)}(T_2)T_{2,t}&=k_2T_{2,xx},\\
\rho_3c_3T_{3,t}&=k_3T_{3,xx}.
\end{aligned}
\right.
\label{eq:q4-governing}
\end{equation}
人体侧边界为
\begin{equation}
-k_1T_{1,x}(0,t)=3.0\,[37-T_1(0,t)],
\label{eq:q4-inner-boundary}
\end{equation}
服装外表面仍采用问题一的无风自然对流关系
\begin{equation}
-k_3T_{3,x}(L,t)
=2.38\,[T_3(L,t)+40]^{1.25}.
\label{eq:q4-outer-boundary}
\end{equation}
两处材料界面继续满足温度和热流连续，三层初始温度均取$37\,^{\circ}\mathrm C$。

定义贴近人体一侧的评价温度
\begin{equation}
T_{\rm in}^{(\alpha)}(t)=T_1(0,t;\alpha),
\end{equation}
并以其首次降至$15\,^{\circ}\mathrm C$的时刻作为放热倍率$\alpha$下的热安全时间
\begin{equation}
t_T^{(4)}(\alpha)
=\inf\{t>0:T_{\rm in}^{(\alpha)}(t)\le15\,^{\circ}\mathrm C\}.
\label{eq:q4-thermal-time}
\end{equation}

问题四继续沿用问题三最终确定的承重口径。由于本问只改变功能材料的单位质量放热能力，不改变服装厚度和质量，因此负重约束对应的最长时间记为常数$t_W^{(4)}$，与$\alpha$无关。实际坚持时间为
\begin{equation}
t_{\rm eff}^{(4)}(\alpha)
=\min\{t_T^{(4)}(\alpha),\ t_W^{(4)}\}.
\label{eq:q4-effective-time}
\end{equation}

\subsection{最小放热能力提高比例模型}

以问题三最优方案得到的实际坚持时间$t_3^*$为性能基准，问题四写成单变量约束优化模型。

\noindent\textbf{1. 决策变量}\par
设$\alpha\ge1$为中间功能层相对于原材料的放热能力倍率，其小数部分$\alpha-1$即为所求提高比例。

\noindent\textbf{2. 目标函数}\par
最小化放热能力提高比例，即
\begin{equation}
\min_{\alpha}\quad \alpha-1.
\label{eq:q4-obj}
\end{equation}

\noindent\textbf{3. 约束条件}\par

\noindent\textbf{（1）坚持时间约束：}通过放热倍率获得的实际坚持时间不得低于问题三的最优水平，即
\begin{equation}
t_{\rm eff}^{(4)}(\alpha)\ge t_3^*.
\label{eq:q4-time-constraint}
\end{equation}

\noindent\textbf{（2）可行性约束：}放热倍率的物理下限为$\alpha\ge1$，对应不放热能力无提高。

\subsubsection{模型汇总}
{\boldmath
\begin{equation}
\min_{\alpha\ge1}\quad \alpha-1
\label{eq:q4-summary-obj}
\end{equation}}

{\boldmath
\begin{equation}
\mathrm{s.t.}\quad
\left\{
\begin{aligned}
&t_{\rm eff}^{(4)}(\alpha)\ge t_3^*,\\
&\alpha\ge1.
\end{aligned}
\right.
\label{eq:q4-summary-con}
\end{equation}}

\subsection{二分搜索求解}

设计变量仅有一个且$t_T^{(4)}(\alpha)$具有明确的单调性，本问无需使用高维启发式方法。先令下界$\alpha_L=1$，逐步扩大上界$\alpha_U$直至$t_T^{(4)}(\alpha_U)\ge t_3^*$；随后对区间$[\alpha_L,\alpha_U]$进行二分。每次取
\begin{equation}
\alpha_M=\frac{\alpha_L+\alpha_U}{2},
\end{equation}
调用问题一的短时间分离变量正演求得$t_T^{(4)}(\alpha_M)$。若$t_T^{(4)}(\alpha_M)\ge t_3^*$，则令$\alpha_U=\alpha_M$；否则令$\alpha_L=\alpha_M$。当倍率区间或坚持时间误差达到预设容差后，取$\alpha^*=\alpha_U$。

在每个微小时间段内，仅需将问题一第二层的有效比热更新为
\begin{equation}
c_{2,\rm eff}^{n,(\alpha)}
=c_2+\alpha L_{\rm pcm}w(\bar T_2^n),
\end{equation}
并取
\begin{equation}
\alpha_2^{n,(\alpha)}
=\frac{k_2}{\rho_2c_{2,\rm eff}^{n,(\alpha)}}.
\end{equation}
第一层、第三层的本征展开、两界面斜率闭合以及外侧换热系数更新时间推进均直接复用问题一，因此问题四只是在同一正演器外增加一个单变量二分搜索层。完整求解步骤整理为算法\ref{alg:q4-bisection}。

\begin{algorithm}[H]
\SetAlgoLined
\setstretch{0.92}
\SetAlFnt{\normalsize}
\SetArgSty{textnormal}
\SetKwInput{KwIn}{输入}
\SetKwInput{KwOut}{输出}
\SetKwFor{For}{for}{do}{end for}
\SetKwIF{If}{ElseIf}{Else}{if}{then}{else if}{else}{end if}
\SetKw{KwRet}{return}
\caption{问题四最小放热倍率二分搜索}
\label{alg:q4-bisection}
\KwIn{基准时间$t_3^*$，负重上限$t_W^{(4)}$，容差$\varepsilon_\alpha$，各扫描速率下潜热$L_{\rm pcm}(v)$}
\KwOut{最小放热倍率$\alpha^*$，提高比例$r_{\min}=100(\alpha^*-1)\%$}

\If{$t_W^{(4)}<t_3^*$}{
  \KwRet 负重上限不足，本问无可行解\;
}
$\alpha_L\leftarrow1$\;
$\alpha_U\leftarrow2$\;
\While{$t_T^{(4)}(\alpha_U)<t_3^*$}{
  $\alpha_L\leftarrow\alpha_U$\;
  $\alpha_U\leftarrow2\alpha_U$\;
}
\While{$\alpha_U-\alpha_L>\varepsilon_\alpha$}{
  $\alpha_M\leftarrow(\alpha_L+\alpha_U)/2$\;
  调用问题一正演器计算$t_T^{(4)}(\alpha_M)$\;
  \If{$t_T^{(4)}(\alpha_M)\ge t_3^*$}{
    $\alpha_U\leftarrow\alpha_M$\;
  }
  \Else{
    $\alpha_L\leftarrow\alpha_M$\;
  }
}
$\alpha^*\leftarrow\alpha_U$\;
$t_{\rm eff}^{(4)}(\alpha^*)\leftarrow\min\{t_T^{(4)}(\alpha^*),t_W^{(4)}\}$\;
\KwRet $\alpha^*$，$100(\alpha^*-1)\%$，$t_{\rm eff}^{(4)}(\alpha^*)$\;
\end{algorithm}

\subsection{求解结果与分析}

\noindent\textbf{1. 求解结果}\par

问题四在$v=2,5,10\ \mathrm{K/min}$三个扫描速率下，各自匹配同一扫描速率下问题一的潜热$L_{\rm pcm}(v)$与问题三的最优实际坚持时间$t_3^*(v)$进行反求。原始厚度下的负重上限均为$t_W^{(4)}=784.916\ \mathrm{s}$，三组均通过可行性检查，且活动约束均为热安全控制。求解结果汇总于表\ref{tab:q4-results}。

\begin{table}[!htbp]
\centering
\caption{问题四各扫描速率下最小放热能力提高比例}
\label{tab:q4-results}
\begin{tabular}{cccccc}
\toprule
\rowcolor{CumcmLightBlue}
$v$\ (\mathrm{K/min}) & $\alpha^*$ & $r_{\min}$(\%) & $t_T^{(4)}$(\mathrm{s}) & $t_{\rm eff}^{(4)}$(\mathrm{s}) & 活动约束 \\
\midrule
2    & 1.437866 & 43.787 & 701.204 & 701.204 & 热安全 \\
5    & 1.863281 & 86.328 & 416.043 & 416.043 & 热安全 \\
10   & 2.572876 & 157.288 & 321.061 & 321.061 & 热安全 \\
\bottomrule
\end{tabular}
\end{table}

\noindent\textbf{2. 结果分析}\par

从表\ref{tab:q4-results}可以看出，扫描速率越低，所需放热能力提高比例越小。$v=2\ \mathrm{K/min}$时仅需提高$43.787\%$，而$v=10\ \mathrm{K/min}$时需提高$157.288\%$。这一趋势的原因在于：低速扫描下相变过程的持续时间更长，潜热$L_{\rm pcm}(v)$更大，原始材料已能较好地延缓温度下降，因此只需较小的倍率提升即可达到目标时间。

三组二分搜索的区间宽度均为$6.104\times10^{-5}$，对应的百分比区间宽度均小于$0.01$个百分点，数值精度满足要求。单调性验证表明，在$\alpha=1,1.5,\alpha^*,2$及二分下界处，热安全时间$t_T^{(4)}(\alpha)$随$\alpha$严格单调递增，未发现平台或异常反向变化，验证结果见表\ref{tab:q4-validation}。

\begin{table}[!htbp]
\centering
\caption{问题四放热倍率单调性验证}
\label{tab:q4-validation}
\begin{tabular}{cccc}
\toprule
\rowcolor{CumcmLightBlue}
$\alpha$ & $v=2$ & $v=5$ & $v=10$ \\
\midrule
1.0      & 520.903 & 273.908 & 191.642 \\
1.5      & 726.790 & 356.223 & 232.766 \\
$\alpha^*$ & 701.204 & 416.043 & 321.061 \\
2.0      & 932.696 & 438.557 & 273.908 \\
\bottomrule
\end{tabular}
\end{table}

需要指出的是，问题三的目标阈值和扫描速率可信域尚未最终锁定，上述结果对应于$15\,^{\circ}\mathrm C$阈值口径下的阶段结果。若问题三最终目标定义调整，问题四需严格继承同一目标定义重新反求。